\newcommand{\ModuleID}{EL-I.12}
\newcommand{\ModuleTitle}{A Repeatable Method for Reading the Missal}
\newcommand{\ModuleStage}{Stage I — Foundations}
\newcommand{\ModuleUnit}{I.12}
\newcommand{\ModuleOutcome}{Apply a seven-pass, evidence-first process to short unfamiliar Missal selections and produce close and natural translations}
\newcommand{\ModulePrerequisites}{EL-I.11 mastery: resolve foundational pronouns and preposition government}
\newcommand{\ModuleSessionOne}{Complete one guided seven-pass reading with a fully annotated scaffold, then complete Worksheet I.45 as the guided cumulative application.}
\newcommand{\ModuleSessionTwo}{Repeat the method with a faded scaffold and disciplined lookup/error records, complete Worksheet I.48 independently, and make the Stage I advancement decision from evidence.}
\newcommand{\ModulePrintedPractice}{Worksheets I.45 (guided Variant A) and I.48 (independent Variant D), with terminal keys}
\newcommand{\ModuleExtensionPractice}{Worksheets I.46 and I.47 remain optional remediation or delayed-recall variants}
\newcommand{\ModuleNext}{EL-GATE-I, Foundations Assessment Sequence (F-A and delayed F-B); after passage, EL-II.1, Relative, Interrogative, and Indefinite Pronouns}
\newcommand{\ModuleMastery}{Within the stated gate conditions, reach at least 90 percent morphological and 85 percent syntactic accuracy, translate every clause without omission, meet the cumulative transformation threshold, and write the requiring rule beside every correction.}
\newcommand{\ModuleDelayNote}{\par\medskip\noindent\textbf{Stage edge.} Advance to Stage II only after the complete Volume I standard is met and the separate EL-GATE-I sequence is passed; elapsed sessions do not waive any threshold.}
\newcommand{\ModuleReferenceFocus}{%
  \begin{itemize}
    \item seven passes: boundaries, finite verbs, nominal forms, agreement, government/reference, close sense, and natural rendering;
    \item printed versus supplied words and visibly bracketed ellipsis;
    \item dictionary use limited to lexical meaning and principal parts after an unaided hypothesis; and
    \item an error ledger that records form, failed decision, corrected rule, and delayed retest.
  \end{itemize}}
\newcommand{\ModuleContrast}{A close translation exposes grammatical relations; natural English communicates the proposition. Neither replaces the morphological and syntactic map, and neither may silently omit a printed word.}
\newcommand{\ModuleLessonSource}{\CourseRoot/shared/content/volume1/all}
\newcommand{\ModuleExerciseSource}{\CourseRoot/shared/exercises/volume1/all}
\newcommand{\ModuleWorkedBank}{I}
\newcommand{\ModuleExerciseBank}{I}
\newcommand{\ModuleWorksheetVariants}{A,D}
\newcommand{\ModuleScopeNote}{This Stage I capstone samples short connected reading; it does not yet establish independent reading of complete unfamiliar formularies.}
\newcommand{\ModuleEnrichment}{%
  \SubmoduleExtension
    {The seven-pass reading method}
    {An annotated-to-faded scaffold makes the method retrievable without turning the scaffold into permanent dependence.}
    {Apply all seven passes to the first guided lesson reading, writing the pass number beside every annotation. Then repeat on another lesson reading with only the seven pass names visible.}
    {List the passes from memory and identify the earliest pass at which an English translation may responsibly be drafted.}
    {The first analysis must account for boundaries, finite verbs, nominal morphology, agreement, government/reference, close sense, and natural rendering. Translation follows the formal map rather than preceding it; wording may be drafted at the close-sense pass and refined only after every earlier decision is explicit.}
  \SubmoduleExtension
    {Guided Missal readings}
    {Fading help tests transfer while retaining the same evidence standard.}
    {Choose one lesson reading whose worked analysis is complete. Cover one layer at a time—translation, syntax, then morphology—and reconstruct each layer before checking. Record the first layer that fails.}
    {Explain why remembering a familiar English text cannot count as a successful direct reading.}
    {A successful reconstruction supplies morphology and syntax from the Latin and then produces a complete close sense. The error record must name the first failed decision. Familiar English may check sense after analysis, but cannot establish forms or relations.}
  \SubmoduleExtension
    {Handwritten synthesis and controlled composition}
    {A lookup/error record distinguishes lexical help from outsourcing the parse.}
    {For the lesson synthesis passage, write an unaided lemma and parse hypothesis before each lookup. After checking, record whether the error was lexical identity, morphology, syntax, reference, or English expression and state the correcting rule.}
    {State the two kinds of information this module permits a dictionary to supply at the gate.}
    {The record must preserve the pre-lookup hypothesis and classify each correction. At the gate, the dictionary supplies lexical meaning and principal parts; it does not supply a prepared translation or replace contextual parsing.}
  \SubmoduleExtension
    {Volume I completion standard}
    {An independent gate must measure current transfer, not page completion or recognition of rehearsed answers.}
    {Complete the lesson's three-passage protocol under its stated conditions. Score morphology, syntax, completeness, and transformations separately; then assign every error to a rule and choose advance or remediation.}
    {Without notes, state all threshold categories and explain why a high average cannot erase a repeated structural error.}
    {Advance requires at least 90 percent morphology, 85 percent syntax, complete clause translation, and the stated cumulative transformation score, followed by rule-based correction. A repeated structural error remains evidence that its dependency is unstable even if other categories raise the average.}
}
